\begin{minipage}[t][\PosterLowerRowHeight][t]{\PosterHalfWidth}\vspace{0pt}
  \begin{PosterPanel}{\PosterLowerRowHeight}{8}{DOES IT TRANSFER? AND WHAT WE DO NOT CLAIM}
    \hyphenpenalty=10000\exhyphenpenalty=10000

    \noindent\begin{minipage}[t]{150mm}\vspace{0pt}
      {\PosterCondensed\bfseries\fontsize{74pt}{78pt}\selectfont\color{IASDarkGreen}5 of 5\par}
      {\FigureSize events pass on a second base design (85\,\% inverter share) with the rule applied unchanged: 0 roots beyond the margin.\par}
    \end{minipage}\hfill%
    \begin{minipage}[t]{262mm}\vspace{0pt}
      {\FigureSize \textbf{Nine-bus bank} (assumed dynamic data): transport on its three PLLs gives $-0.131\ \mathrm{s^{-1}}$ from $+0.886$ (spectral re-run; its 12 of 12 events are reported).\par}\vspace{2mm}
      {\FigureSize \textbf{First-order version} of the formula gives the same root counts: the choice of mode matters, not the truncation.\par}\vspace{2mm}
      {\FigureSize \textbf{Phase budget} $h^{\star}=\varphi_{PI}(\omega)/\omega$: 8.99 ms for the 4.91 Hz mode, a limit of the formula, not of operation.\par}
    \end{minipage}\par
    \vspace{4mm}
    \begin{InWords}\InWordsLabel \textbf{Not claimed:} an optimal mode, a maximum delay, a maximum inverter share, superiority over other compensators. Root counts are floating point, checked against eigenvalues at zero delay, not interval-certified. One network.\end{InWords}

  \end{PosterPanel}
\end{minipage}%
